%--- intro 1.2 (reescrita) ---
\subsection{Comparativa de 5 dispositivos móviles}

Para poder comparar algo primero hay que decidir qué se compara. He cogido cinco teléfonos que sean intencionadamente distintos entre sí: tres de gama alta Android y de fabricantes distintos (Samsung, Google y OnePlus), un iPhone de gama alta, que además es el único iOS de la lista y me sirve para ver las diferencias de plataforma, y un Motorola de gama media-baja, que es el que más se parece a lo que usa la gente que no se compra un móvil cada dos años. Los datos salen de la ficha técnica de cada fabricante, que dejo enlazada al final del apartado.

%--- tablas 1.2.1 ---
\subsubsection*{Tabla 1.2.1 --- Comparativa general}

\begin{xltabular}{\textwidth}{@{}>{\RaggedRight\arraybackslash}p{2.4cm} >{\RaggedRight\arraybackslash}X >{\RaggedRight\arraybackslash}X >{\RaggedRight\arraybackslash}X >{\RaggedRight\arraybackslash}X >{\RaggedRight\arraybackslash}p{1.3cm}@{}}
\caption{Comparativa general de los cinco dispositivos}\label{tab:comparativa-general}\\
\toprule
\rowcolor{primary}
\textcolor{white}{\sffamily\bfseries Dispositivo} & \textcolor{white}{\sffamily\bfseries Procesamiento} & \textcolor{white}{\sffamily\bfseries Tamaño} & \textcolor{white}{\sffamily\bfseries Autonomía} & \textcolor{white}{\sffamily\bfseries Conectividad} & \textcolor{white}{\sffamily\bfseries Fuente} \\
\midrule
\endfirsthead
\toprule
\rowcolor{primary}
\textcolor{white}{\sffamily\bfseries Dispositivo} & \textcolor{white}{\sffamily\bfseries Procesamiento} & \textcolor{white}{\sffamily\bfseries Tamaño} & \textcolor{white}{\sffamily\bfseries Autonomía} & \textcolor{white}{\sffamily\bfseries Conectividad} & \textcolor{white}{\sffamily\bfseries Fuente} \\
\midrule
\endhead
\bottomrule
\endlastfoot
\textbf{Samsung Galaxy S26 Ultra} &
Snapdragon 8 Elite Gen 5, 2$\times$4,74 + 6$\times$3,62 GHz · 12-16 GB RAM · 256 GB-1 TB &
6,9'' · 1440$\times$3120 · Dynamic AMOLED 2X · 163,6$\times$78,1$\times$7,9 mm · 214 g &
5.000 mAh · 60 W · hasta 31 h de vídeo &
5G, Wi-Fi 7, BT 6.0, NFC, UWB, USB-C 3.2, satélite &
\href{https://www.samsung.com/es/smartphones/galaxy-s26-ultra/}{Samsung} \\
\rowcolor{soft}
\textbf{Google Pixel 10 Pro XL} &
Tensor G5, 1$\times$3,78 + 5$\times$3,05 + 2$\times$2,25 GHz · 16 GB RAM · 256 GB-1 TB &
6,8'' · 1344$\times$2992 · LTPO OLED · 162,8$\times$76,6$\times$8,5 mm · 232 g &
5.200 mAh · 45 W · más de 24 h &
5G, Wi-Fi 7, BT 6.0, NFC, UWB, Thread, USB-C 3.2, satélite &
\href{https://store.google.com/product/pixel_10_pro_specs}{Google} \\
\textbf{OnePlus 13} &
Snapdragon 8 Elite, 2$\times$4,32 + 6$\times$3,53 GHz · 12-16 GB RAM · 256-512 GB &
6,82'' · 1440$\times$3168 · LTPO ProXDR 120 Hz · 162,9$\times$76,5$\times$8,5 mm · 213 g &
6.000 mAh · 80 W (100 W opcional) · sin cifra oficial &
5G, Wi-Fi 7, BT 5,4, NFC, mando IR, USB 3.2 &
\href{https://www.oneplus.com/es/13/specs}{OnePlus} \\
\rowcolor{soft}
\textbf{iPhone 17 Pro Max} &
A19 Pro, 6 núcleos (2 rendimiento + 4 eficiencia) · 12 GB RAM · 256 GB-2 TB &
6,9'' · 1320$\times$2868 · Super Retina XDR OLED · 163,4$\times$78$\times$8,75 mm · 233 g &
$\sim$5.088 mAh (no publicado) · MagSafe 25 W · hasta 39 h de vídeo &
5G, Wi-Fi 7, BT 6, NFC, UWB, Thread, USB-C 10 Gb/s, satélite; solo eSIM en EE.UU. &
\href{https://support.apple.com/es-es/125091}{Apple} \\
\textbf{Motorola Moto G47} &
MediaTek Dimensity 6300, 8 núcleos a 2,4 GHz · 8 GB RAM (+16 GB virtual) · 128 GB + microSD 1 TB &
6,67'' · 720$\times$1604 · LCD · 166,2$\times$76,5$\times$7,85 mm · 191 g &
5.200 mAh · 20 W · hasta 43 h &
5G, Wi-Fi 5, BT 5,4, NFC, jack 3,5 mm, USB-C 2.0 &
\href{https://www.motorola.com/sa/en/p/phones/moto-g/47/pmipmjs44mz}{Motorola} \\
\end{xltabular}

\begin{warnbox}
\textbf{\color{accent!80!black}\faExclamationTriangle\ \ Dos avisos sobre los datos.}
De todo lo que he ido metiendo en las tablas hay dos cosas que no son reales y conviene que se
lean bien. La primera es que \textbf{Apple no publica ni los mAh de la batería ni los GB de RAM del
iPhone 17 Pro Max}; esas dos cifras vienen de desmontajes de terceros y las he marcado como tales en
la tabla. La segunda es el \textbf{RAM Boost del Moto G47}, que no es RAM real: Motorola usa parte del
almacenamiento interno como memoria virtual para anunciar hasta 24 GB. Sirve para algo, pero no es lo
mismo que 24 GB de RAM, y es un error bastante habitual.
\end{warnbox}
%--- tablas 1.2.2 y 1.2.3 ---
\subsubsection*{Tabla 1.2.2 --- A. Capacidad de procesamiento}

\begin{xltabular}{\textwidth}{@{}>{\RaggedRight\arraybackslash}p{2.3cm} >{\RaggedRight\arraybackslash}X >{\RaggedRight\arraybackslash}X >{\RaggedRight\arraybackslash}X >{\RaggedRight\arraybackslash}X >{\RaggedRight\arraybackslash}X@{}}
\caption{Comparativa de capacidad de procesamiento}\label{tab:proc}\\
\toprule
\rowcolor{primary}
\textcolor{white}{\sffamily\bfseries} & \textcolor{white}{\sffamily\bfseries Samsung S26 Ultra} & \textcolor{white}{\sffamily\bfseries Pixel 10 Pro XL} & \textcolor{white}{\sffamily\bfseries OnePlus 13} & \textcolor{white}{\sffamily\bfseries iPhone 17 Pro Max} & \textcolor{white}{\sffamily\bfseries Moto G47} \\
\midrule
\endfirsthead
\toprule
\rowcolor{primary}
\textcolor{white}{\sffamily\bfseries} & \textcolor{white}{\sffamily\bfseries Samsung} & \textcolor{white}{\sffamily\bfseries Pixel 10 Pro XL} & \textcolor{white}{\sffamily\bfseries OnePlus 13} & \textcolor{white}{\sffamily\bfseries iPhone 17 Pro Max} & \textcolor{white}{\sffamily\bfseries Moto G47} \\
\midrule
\endhead
\bottomrule
\endlastfoot
\textbf{Procesador / SoC} &
Qualcomm Snapdragon 8 Elite Gen 5 for Galaxy (TSMC 3 nm) &
Google Tensor G5 (TSMC 3 nm) + coprocesador Titan M2 &
Qualcomm Snapdragon 8 Elite (TSMC 3 nm), módem X80 &
Apple A19 Pro (TSMC 3 nm, según terceros) &
MediaTek Dimensity 6300 (6 nm) \\
\rowcolor{soft}
\textbf{Núcleos} &
Octa-core: 2$\times$ Oryon V3 a 4,74 GHz + 6$\times$ Oryon V3 a 3,62 GHz &
Octa-core: 1$\times$ Cortex-X4 3,78 GHz + 5$\times$ Cortex-A725 3,05 GHz + 2$\times$ A520 2,25 GHz &
Octa-core: 2$\times$ Oryon a 4,32 GHz + 6$\times$ Oryon a 3,53 GHz &
Hexa-core: 2 rendimiento + 4 eficiencia &
Octa-core: 2$\times$ A76 2,4 GHz + 6$\times$ A55 2,0 GHz \\
\textbf{GPU} & Adreno 840 & PowerVR DXT-48-1536 & Adreno 830 & 6 núcleos con ray tracing & Arm Mali-G57 M2C \\
\rowcolor{soft}
\textbf{RAM} &
12 GB (256/512 GB) · 16 GB (1 TB) &
16 GB en todas las versiones &
12 / 16 GB LPDDR5X &
12 GB (no publicado) &
8 GB LPDDR4X (+16 GB virtuales) \\
\textbf{Almacenamiento} & 256 GB / 512 GB / 1 TB (UFS 4.x) & 256 GB / 512 GB / 1 TB (UFS 4.0) & 256 GB / 512 GB (UFS 4.0) & 256 GB / 512 GB / 1 TB / 2 TB & 128 GB + microSD hasta 1 TB \\
\rowcolor{soft}
\textbf{Ampliable} & No & No & No & No & \textbf{Sí (único de los cinco)} \\
\end{xltabular}

\subsubsection*{Tabla 1.2.3 --- B. Tamaño}

\begin{xltabular}{\textwidth}{@{}>{\RaggedRight\arraybackslash}p{2.3cm} >{\RaggedRight\arraybackslash}X >{\RaggedRight\arraybackslash}X >{\RaggedRight\arraybackslash}X >{\RaggedRight\arraybackslash}X >{\RaggedRight\arraybackslash}X@{}}
\caption{Comparativa de tamaño y pantalla}\label{tab:tamano}\\
\toprule
\rowcolor{primary}
\textcolor{white}{\sffamily\bfseries} & \textcolor{white}{\sffamily\bfseries Samsung} & \textcolor{white}{\sffamily\bfseries Pixel 10 Pro XL} & \textcolor{white}{\sffamily\bfseries OnePlus 13} & \textcolor{white}{\sffamily\bfseries iPhone 17 Pro Max} & \textcolor{white}{\sffamily\bfseries Moto G47} \\
\midrule
\endfirsthead
\toprule
\rowcolor{primary}
\textcolor{white}{\sffamily\bfseries} & \textcolor{white}{\sffamily\bfseries Samsung} & \textcolor{white}{\sffamily\bfseries Pixel 10 Pro XL} & \textcolor{white}{\sffamily\bfseries OnePlus 13} & \textcolor{white}{\sffamily\bfseries iPhone 17 Pro Max} & \textcolor{white}{\sffamily\bfseries Moto G47} \\
\midrule
\endhead
\bottomrule
\endlastfoot
\textbf{Pulgadas} & 6,9'' & 6,8'' & 6,82'' & 6,9'' & 6,67'' \\
\rowcolor{soft}
\textbf{Resolución} & 1440$\times$3120 (QHD+) & 1344$\times$2992 (QHD+) & 1440$\times$3168 (QHD+) & 1320$\times$2868 & \textbf{720$\times$1604 (HD+)} \\
\textbf{Densidad aprox.} & $\sim$500 ppi & $\sim$486 ppi & $\sim$510 ppi & 460 ppi & \textbf{264 ppi} \\
\rowcolor{soft}
\textbf{Tecnología} & Dynamic AMOLED 2X (LTPO), 2.600 nits, Gorilla Armor 2 & LTPO OLED, 3.300 nits, Victus 2 & ProXDR LTPO 4.1, 4.500 nits, 10 bits & Super Retina XDR OLED, 3.000 nits & \textbf{LCD de 8 bits}, 1.050 nits \\
\textbf{Refresco} & 120 Hz adaptativo & 1-120 Hz & 1-120 Hz dinámico & ProMotion hasta 120 Hz & 120 Hz \\
\rowcolor{soft}
\textbf{Dimensiones} & 163,6$\times$78,1$\times$7,9 mm & 162,8$\times$76,6$\times$8,5 mm & 162,9$\times$76,5$\times$8,5 mm & 163,4$\times$78,0$\times$8,75 mm & 166,2$\times$76,5$\times$7,85 mm \\
\textbf{Peso} & 214 g & 232 g & 213 g & \textbf{233 g (el más pesado)} & \textbf{191 g (el más ligero)} \\
\end{xltabular}
%--- tablas 1.2.4, 1.2.5 y fuentes (se omite la caja "Fjate en...") ---
\subsubsection*{Tabla 1.2.4 --- C. Autonomía}

\begin{xltabular}{\textwidth}{@{}>{\RaggedRight\arraybackslash}p{2.7cm} >{\RaggedRight\arraybackslash}X >{\RaggedRight\arraybackslash}X >{\RaggedRight\arraybackslash}X >{\RaggedRight\arraybackslash}X >{\RaggedRight\arraybackslash}X@{}}
\caption{Comparativa de autonomía y carga}\label{tab:bateria}\\
\toprule
\rowcolor{primary}
\textcolor{white}{\sffamily\bfseries} & \textcolor{white}{\sffamily\bfseries Samsung} & \textcolor{white}{\sffamily\bfseries Pixel 10 Pro XL} & \textcolor{white}{\sffamily\bfseries OnePlus 13} & \textcolor{white}{\sffamily\bfseries iPhone 17 Pro Max} & \textcolor{white}{\sffamily\bfseries Moto G47} \\
\midrule
\endfirsthead
\toprule
\rowcolor{primary}
\textcolor{white}{\sffamily\bfseries} & \textcolor{white}{\sffamily\bfseries Samsung} & \textcolor{white}{\sffamily\bfseries Pixel 10 Pro XL} & \textcolor{white}{\sffamily\bfseries OnePlus 13} & \textcolor{white}{\sffamily\bfseries iPhone 17 Pro Max} & \textcolor{white}{\sffamily\bfseries Moto G47} \\
\midrule
\endhead
\bottomrule
\endlastfoot
\textbf{Batería} & 5.000 mAh & 5.200 mAh & \textbf{6.000 mAh (el más grande)} & $\sim$5.088 mAh (terceros) & 5.200 mAh \\
\rowcolor{soft}
\textbf{Carga con cable} & 60 W (75\,\% en 30 min) & 45 W (70\,\% en 30 min) & \textbf{80 W, hasta 100 W con adaptador} & 50\,\% en 20 min con 40 W & 20 W \\
\textbf{Carga inalámbrica} & 25 W & 25 W (con imanes) & \textbf{50 W} & MagSafe 25 W / Qi2 25 W & \textbf{No tiene} \\
\rowcolor{soft}
\textbf{Declaración} & Hasta \textbf{31 h} de vídeo & Más de \textbf{24 h} & \textbf{No publica horas} & Hasta \textbf{39 h} de vídeo & Hasta \textbf{43 h} (varía) \\
\textbf{Extra} & Carga inversa 4,5 W & Carga inversa 10 W & Carga inversa 10 W & Sin cargador en la caja & Carga inversa 6 W \\
\end{xltabular}

\subsubsection*{Tabla 1.2.5 --- D. Conectividad}

\begin{xltabular}{\textwidth}{@{}>{\RaggedRight\arraybackslash}p{2.4cm} >{\RaggedRight\arraybackslash}X >{\RaggedRight\arraybackslash}X >{\RaggedRight\arraybackslash}X >{\RaggedRight\arraybackslash}X >{\RaggedRight\arraybackslash}X@{}}
\caption{Comparativa de conectividad}\label{tab:conectividad}\\
\toprule
\rowcolor{primary}
\textcolor{white}{\sffamily\bfseries} & \textcolor{white}{\sffamily\bfseries Samsung} & \textcolor{white}{\sffamily\bfseries Pixel 10 Pro XL} & \textcolor{white}{\sffamily\bfseries OnePlus 13} & \textcolor{white}{\sffamily\bfseries iPhone 17 Pro Max} & \textcolor{white}{\sffamily\bfseries Moto G47} \\
\midrule
\endfirsthead
\toprule
\rowcolor{primary}
\textcolor{white}{\sffamily\bfseries} & \textcolor{white}{\sffamily\bfseries Samsung} & \textcolor{white}{\sffamily\bfseries Pixel 10 Pro XL} & \textcolor{white}{\sffamily\bfseries OnePlus 13} & \textcolor{white}{\sffamily\bfseries iPhone 17 Pro Max} & \textcolor{white}{\sffamily\bfseries Moto G47} \\
\midrule
\endhead
\bottomrule
\endlastfoot
\textbf{Redes móviles} & 2G/3G/4G/5G (SA y NSA) + mmWave EE.UU. & 2G/3G/4G/5G (SA y NSA) & 2G/3G/4G/5G (solo Sub-6) & 5G sub-6 + mmWave, LTE gigabit & 2G/3G/4G/5G (solo Sub-6) \\
\rowcolor{soft}
\textbf{Wi-Fi} & Wi-Fi 7 (802.11be) & Wi-Fi 7 (802.11be) & Wi-Fi 7 (802.11be) & Wi-Fi 7 (802.11be) & \textbf{Wi-Fi 5 (802.11ac)} \\
\textbf{Bluetooth} & \textbf{6.0} & \textbf{6.0} & 5,4 & \textbf{6} & 5,4 \\
\rowcolor{soft}
\textbf{NFC} & Sí (con UWB) & Sí & Sí (con eSE) & Sí (modo lector) & Sí (según mercado) \\
\textbf{UWB} & Sí & Sí & \textbf{No} & Sí (chip 2.ª gen.) & \textbf{No} \\
\rowcolor{soft}
\textbf{Otros detalles} & Satélite (SOS y mensajes) & Satélite SOS, Thread/Matter & \textbf{Mando IR}, BeaconLink & Satélite, Thread, USB-C 10 Gb/s & \textbf{Jack 3,5 mm}, IP64, MIL-STD-810H \\
\textbf{SIM} & Nano-SIM + eSIM & Nano-SIM + eSIM & Doble nano-SIM & \textbf{Solo eSIM EE.UU.}, nano+eSIM Europa & 2$\times$ nano-SIM, \textbf{sin eSIM} \\
\rowcolor{soft}
\textbf{USB} & USB-C 3.2 & USB-C 3.2 & USB 3.2 & USB-C hasta 10 Gb/s & USB-C 2.0 \\
\end{xltabular}

\begin{fuentesbox}
\textbf{\color{muted}\faBook\ \ Fuentes del apartado 1.2:}
\href{https://www.samsung.com/es/smartphones/galaxy-s26-ultra/}{Samsung Galaxy S26 Ultra} ·
\href{https://store.google.com/product/pixel_10_pro_specs}{Google Pixel 10 Pro XL} ·
\href{https://www.oneplus.com/es/13/specs}{OnePlus 13} ·
\href{https://support.apple.com/es-es/125091}{iPhone 17 Pro Max} ·
\href{https://www.motorola.com/sa/en/p/phones/moto-g/47/pmipmjs44mz}{Motorola Moto G47} ·
Secundaria: \href{https://www.gsmarena.com/}{GSMArena}
\end{fuentesbox}
